% ============================================================================
% CHAPTER 7: ADVANCED TOPICS
% ============================================================================

\chapter{Advanced Topics}

\section{Custom Model Training}

\subsection{Training Custom BERT Classifier}

For domain-specific performance improvements, you can train a custom BERT model.

\textbf{Prepare Training Data:}
\begin{lstlisting}[style=pythonstyle]
# Format: CSV with columns: text, label
# ml_training/data/training_data.csv

text,label
"55-year-old male developed rash after aspirin",ICSR
"Clinical trial showed 30% response rate",Safety_Information
"Aspirin inhibits COX-2 enzyme",Non-ICSR
"FDA approves new cancer drug",Rejected
\end{lstlisting}

\textbf{Training Script:}
\begin{lstlisting}[style=pythonstyle]
# ml_training/train_bert.py
from transformers import AutoModelForSequenceClassification, AutoTokenizer
from transformers import Trainer, TrainingArguments
import pandas as pd

# Load data
df = pd.read_csv('data/training_data.csv')
texts = df['text'].tolist()
labels = df['label'].map({
    'ICSR': 0,
    'Safety_Information': 1,
    'Non-ICSR': 2,
    'Rejected': 3
}).tolist()

# Load pre-trained model
model_name = "facebook/bart-large-mnli"
model = AutoModelForSequenceClassification.from_pretrained(
    model_name,
    num_labels=4
)
tokenizer = AutoTokenizer.from_pretrained(model_name)

# Training arguments
training_args = TrainingArguments(
    output_dir='./results',
    num_train_epochs=3,
    per_device_train_batch_size=8,
    save_steps=500,
    logging_steps=100,
    evaluation_strategy="steps",
    learning_rate=2e-5
)

# Train
trainer = Trainer(
    model=model,
    args=training_args,
    train_dataset=train_dataset,
    eval_dataset=eval_dataset
)
trainer.train()

# Save model
model.save_pretrained('./custom_bert_model')
tokenizer.save_pretrained('./custom_bert_model')
\end{lstlisting}

\textbf{Deploy Custom Model:}
\begin{lstlisting}[style=pythonstyle]
# Update services/pv-classification-service/bert_classifier.py
MODEL_NAME = "/app/models/custom_bert_model"

# Copy model to container
# docker-compose.yml
volumes:
  - ./ml_training/custom_bert_model:/app/models/custom_bert_model
\end{lstlisting}

\subsection{Reinforcement Learning Model Retraining}

\textbf{Collect New Training Data:}
\begin{commandbox}
\begin{lstlisting}[style=bashstyle]
# Collect 500 recent documents with classification results
python rl_optimization/data_collection/collect_real_pubmed_data.py \
  --max-results 500 \
  --output rl_optimization/real_pubmed_data
\end{lstlisting}
\end{commandbox}

\textbf{Retrain RL Agent:}
\begin{commandbox}
\begin{lstlisting}[style=bashstyle]
# Train with new data
python rl_optimization/rl_training/training_engine.py \
  --episodes 1000 \
  --data-path rl_optimization/real_pubmed_data \
  --checkpoint-dir rl_optimization/checkpoints

# Results saved to: rl_optimization/training_results/
\end{lstlisting}
\end{commandbox}

\textbf{Deploy Updated Model:}
\begin{commandbox}
\begin{lstlisting}[style=bashstyle]
# Copy new model to RL Intelligence Service
docker-compose exec rl-intelligence-service cp \
  /app/rl_optimization/checkpoints/best_model.pt \
  /app/models/rl_agent.pt

# Restart service
docker-compose restart rl-intelligence-service
\end{lstlisting}
\end{commandbox}

\section{API Integration Examples}

\subsection{Python SDK Example}

\begin{lstlisting}[style=pythonstyle]
# example_client.py
import requests
from typing import List, Dict

class DataIngestionClient:
    def __init__(self, base_url: str, email: str, password: str):
        self.base_url = base_url
        self.token = self._authenticate(email, password)
    
    def _authenticate(self, email: str, password: str) -> str:
        """Authenticate and get access token"""
        response = requests.post(
            f"{self.base_url}/api/v1/auth/login",
            json={"email": email, "password": password}
        )
        response.raise_for_status()
        return response.json()["access_token"]
    
    def _headers(self) -> Dict:
        """Get authorization headers"""
        return {
            "Authorization": f"Bearer {self.token}",
            "Content-Type": "application/json"
        }
    
    def upload_document(self, file_path: str, space_id: str) -> Dict:
        """Upload and process document"""
        with open(file_path, 'rb') as f:
            files = {'files': f}
            data = {
                'space_id': space_id,
                'processing_options': {
                    'enable_classification': True,
                    'enable_entity_extraction': True
                }
            }
            response = requests.post(
                f"{self.base_url}/api/v1/upload",
                headers={"Authorization": f"Bearer {self.token}"},
                files=files,
                data={'processing_options': str(data['processing_options'])}
            )
            response.raise_for_status()
            return response.json()
    
    def get_document(self, doc_id: str) -> Dict:
        """Retrieve document details"""
        response = requests.get(
            f"{self.base_url}/api/v1/documents/{doc_id}",
            headers=self._headers()
        )
        response.raise_for_status()
        return response.json()
    
    def classify_text(self, title: str, abstract: str) -> Dict:
        """Classify document text"""
        response = requests.post(
            f"{self.base_url}/api/v1/pv-classification/classify",
            headers=self._headers(),
            json={
                "title": title,
                "abstract": abstract,
                "force_ensemble": False
            }
        )
        response.raise_for_status()
        return response.json()
    
    def get_knowledge_graph(self, limit: int = 500) -> Dict:
        """Retrieve knowledge graph data"""
        response = requests.get(
            f"{self.base_url}/api/v1/knowledge-hub/graph?limit={limit}",
            headers=self._headers()
        )
        response.raise_for_status()
        return response.json()

# Usage example
if __name__ == "__main__":
    client = DataIngestionClient(
        base_url="http://localhost:8000",
        email="admin@example.com",
        password="admin123"
    )
    
    # Upload document
    result = client.upload_document(
        file_path="./study_report.pdf",
        space_id="space_abc123"
    )
    print(f"Document uploaded: {result['documents'][0]['doc_id']}")
    
    # Classify text
    classification = client.classify_text(
        title="Adverse Event Case Report",
        abstract="Patient developed severe rash after aspirin administration"
    )
    print(f"Classification: {classification['classification']}")
    print(f"Confidence: {classification['classification_confidence']}")
\end{lstlisting}

\subsection{JavaScript/Node.js Example}

\begin{lstlisting}[style=pythonstyle]
// example_client.js
const axios = require('axios');
const FormData = require('form-data');
const fs = require('fs');

class DataIngestionClient {
    constructor(baseURL, email, password) {
        this.baseURL = baseURL;
        this.client = axios.create({ baseURL });
        this.authenticate(email, password);
    }

    async authenticate(email, password) {
        const response = await this.client.post('/api/v1/auth/login', {
            email,
            password
        });
        this.token = response.data.access_token;
        this.client.defaults.headers.common['Authorization'] = `Bearer ${this.token}`;
    }

    async uploadDocument(filePath, spaceId) {
        const form = new FormData();
        form.append('files', fs.createReadStream(filePath));
        form.append('space_id', spaceId);
        form.append('processing_options', JSON.stringify({
            enable_classification: true,
            enable_entity_extraction: true
        }));

        const response = await this.client.post('/api/v1/upload', form, {
            headers: form.getHeaders()
        });
        return response.data;
    }

    async getDocument(docId) {
        const response = await this.client.get(`/api/v1/documents/${docId}`);
        return response.data;
    }

    async classifyText(title, abstract) {
        const response = await this.client.post('/api/v1/pv-classification/classify', {
            title,
            abstract,
            force_ensemble: false
        });
        return response.data;
    }

    async getKnowledgeGraph(limit = 500) {
        const response = await this.client.get(`/api/v1/knowledge-hub/graph?limit=${limit}`);
        return response.data;
    }
}

// Usage
(async () => {
    const client = new DataIngestionClient(
        'http://localhost:8000',
        'admin@example.com',
        'admin123'
    );

    // Classify text
    const result = await client.classifyText(
        'Adverse Event Report',
        'Patient developed rash after aspirin'
    );
    console.log(`Classification: ${result.classification}`);
    console.log(`Confidence: ${result.classification_confidence}`);
})();
\end{lstlisting}

\section{Custom Workflows}

\subsection{Automated Daily PubMed Ingestion}

\textbf{Cron Job Script:}
\begin{lstlisting}[style=bashstyle]
#!/bin/bash
# daily_pubmed_ingestion.sh

API_URL="http://localhost:8000"
TOKEN="<your_long_lived_token>"
SPACE_ID="space_abc123"

# Define queries
QUERIES=(
    "COVID-19 adverse events"
    "vaccine safety pharmacovigilance"
    "drug safety case reports"
)

# Start ingestion for each query
for QUERY in "${QUERIES[@]}"; do
    echo "Starting ingestion for: $QUERY"
    curl -X POST "$API_URL/api/v1/ingestion/pubmed/start" \
        -H "Authorization: Bearer $TOKEN" \
        -H "Content-Type: application/json" \
        -d "{
            \"space_id\": \"$SPACE_ID\",
            \"queries\": [\"$QUERY\"],
            \"max_results_per_query\": 50,
            \"date_range\": {
                \"start\": \"$(date -d '1 day ago' +%Y-%m-%d)\",
                \"end\": \"$(date +%Y-%m-%d)\"
            },
            \"processing_options\": {
                \"enable_classification\": true,
                \"enable_entity_extraction\": true
            }
        }"
    
    # Wait 5 minutes between queries
    sleep 300
done

echo "Daily ingestion completed"
\end{lstlisting}

\textbf{Schedule with Cron:}
\begin{commandbox}
\begin{lstlisting}[style=bashstyle]
# Run daily at 3 AM
0 3 * * * /path/to/daily_pubmed_ingestion.sh >> /var/log/ingestion.log 2>&1
\end{lstlisting}
\end{commandbox}

\subsection{Automated ICSR Export to Regulatory System}

\begin{lstlisting}[style=pythonstyle]
# export_icsrs_to_regulatory.py
import requests
import xml.etree.ElementTree as ET
from datetime import datetime, timedelta

API_URL = "http://localhost:8000"
TOKEN = "<your_token>"

def get_new_icsrs(days=1):
    """Fetch ICSRs approved in last N days"""
    since = (datetime.now() - timedelta(days=days)).isoformat()
    response = requests.get(
        f"{API_URL}/api/v1/documents",
        headers={"Authorization": f"Bearer {TOKEN}"},
        params={
            "classification": "ICSR",
            "approved": True,
            "created_after": since,
            "limit": 100
        }
    )
    return response.json()["documents"]

def convert_to_e2b_xml(icsr):
    """Convert ICSR to ICH E2B XML format"""
    root = ET.Element("ichicsr", lang="en")
    
    # Safety report
    safety_report = ET.SubElement(root, "safetyreport")
    
    # Patient information
    patient = ET.SubElement(safety_report, "patient")
    ET.SubElement(patient, "patientinitial").text = icsr["patient"]["initials"]
    ET.SubElement(patient, "patientage").text = str(icsr["patient"]["age"])
    
    # Drug information
    drug = ET.SubElement(safety_report, "drug")
    ET.SubElement(drug, "medicinalproduct").text = icsr["drug"]["name"]
    ET.SubElement(drug, "drugdosagetext").text = icsr["drug"]["dosage"]
    
    # Adverse event
    reaction = ET.SubElement(safety_report, "reaction")
    ET.SubElement(reaction, "primarysourcereaction").text = icsr["event"]["description"]
    
    return ET.tostring(root, encoding='unicode')

def submit_to_regulatory_system(xml_data):
    """Submit E2B XML to regulatory system (mock)"""
    # Replace with actual regulatory system API
    response = requests.post(
        "https://regulatory-system.gov/api/submit",
        headers={"Content-Type": "application/xml"},
        data=xml_data
    )
    return response.json()

if __name__ == "__main__":
    # Get new ICSRs
    icsrs = get_new_icsrs(days=1)
    print(f"Found {len(icsrs)} new ICSRs")
    
    # Export each ICSR
    for icsr in icsrs:
        xml_data = convert_to_e2b_xml(icsr)
        result = submit_to_regulatory_system(xml_data)
        print(f"Submitted ICSR {icsr['doc_id']}: {result['status']}")
\end{lstlisting}

\section{Deployment Strategies}

\subsection{Production Deployment on AWS}

\textbf{Infrastructure as Code (Terraform):}
\begin{lstlisting}[style=pythonstyle]
# terraform/main.tf
provider "aws" {
  region = "us-east-1"
}

# ECS Cluster
resource "aws_ecs_cluster" "main" {
  name = "data-ingestion-cluster"
}

# RDS PostgreSQL
resource "aws_rds_instance" "postgres" {
  identifier           = "data-ingestion-postgres"
  engine              = "postgres"
  engine_version      = "15.4"
  instance_class      = "db.t3.large"
  allocated_storage   = 100
  username            = "postgres"
  password            = var.db_password
  
  vpc_security_group_ids = [aws_security_group.rds.id]
  db_subnet_group_name   = aws_db_subnet_group.main.name
}

# Elasticsearch (OpenSearch Service)
resource "aws_elasticsearch_domain" "main" {
  domain_name           = "data-ingestion-es"
  elasticsearch_version = "8.0"
  
  cluster_config {
    instance_type  = "r6g.large.elasticsearch"
    instance_count = 3
  }
  
  ebs_options {
    ebs_enabled = true
    volume_size = 100
  }
}

# Application Load Balancer
resource "aws_lb" "main" {
  name               = "data-ingestion-alb"
  internal           = false
  load_balancer_type = "application"
  security_groups    = [aws_security_group.alb.id]
  subnets            = aws_subnet.public[*].id
}

# ECS Task Definitions (services)
resource "aws_ecs_task_definition" "frontend" {
  family                   = "frontend"
  network_mode             = "awsvpc"
  requires_compatibilities = ["FARGATE"]
  cpu                      = "512"
  memory                   = "1024"
  
  container_definitions = jsonencode([{
    name  = "frontend"
    image = "your-ecr-repo/frontend:latest"
    portMappings = [{
      containerPort = 80
      protocol      = "tcp"
    }]
  }])
}

# Similar task definitions for other services...
\end{lstlisting}

\textbf{Deploy:}
\begin{commandbox}
\begin{lstlisting}[style=bashstyle]
# Initialize Terraform
terraform init

# Plan deployment
terraform plan -out=plan.tfplan

# Apply
terraform apply plan.tfplan
\end{lstlisting}
\end{commandbox}

\subsection{Kubernetes Deployment}

\textbf{Kubernetes Manifests:}
\begin{lstlisting}[style=pythonstyle]
# k8s/deployment.yaml
apiVersion: apps/v1
kind: Deployment
metadata:
  name: api-gateway
spec:
  replicas: 3
  selector:
    matchLabels:
      app: api-gateway
  template:
    metadata:
      labels:
        app: api-gateway
    spec:
      containers:
      - name: api-gateway
        image: your-registry/api-gateway:latest
        ports:
        - containerPort: 8000
        env:
        - name: POSTGRES_HOST
          value: postgres-service
        - name: ELASTICSEARCH_URL
          value: http://elasticsearch-service:9200
        resources:
          requests:
            memory: "1Gi"
            cpu: "500m"
          limits:
            memory: "2Gi"
            cpu: "1000m"
---
apiVersion: v1
kind: Service
metadata:
  name: api-gateway-service
spec:
  type: LoadBalancer
  ports:
  - port: 80
    targetPort: 8000
  selector:
    app: api-gateway
\end{lstlisting}

\textbf{Deploy to Kubernetes:}
\begin{commandbox}
\begin{lstlisting}[style=bashstyle]
# Apply all manifests
kubectl apply -f k8s/

# Check deployment status
kubectl get deployments
kubectl get pods
kubectl get services

# Scale deployment
kubectl scale deployment api-gateway --replicas=5

# Update deployment (rolling update)
kubectl set image deployment/api-gateway \
  api-gateway=your-registry/api-gateway:v2.0
\end{lstlisting}
\end{commandbox}

\section{Security Hardening}

\subsection{Enable HTTPS/TLS}

\textbf{Production Nginx Configuration:}
\begin{lstlisting}[style=pythonstyle]
# frontend/nginx-prod.conf
server {
    listen 80;
    server_name yourdomain.com;
    return 301 https://$server_name$request_uri;
}

server {
    listen 443 ssl http2;
    server_name yourdomain.com;
    
    # SSL certificates (Let's Encrypt)
    ssl_certificate /etc/letsencrypt/live/yourdomain.com/fullchain.pem;
    ssl_certificate_key /etc/letsencrypt/live/yourdomain.com/privkey.pem;
    
    # SSL configuration
    ssl_protocols TLSv1.3 TLSv1.2;
    ssl_ciphers 'ECDHE-ECDSA-AES256-GCM-SHA384:ECDHE-RSA-AES256-GCM-SHA384';
    ssl_prefer_server_ciphers on;
    ssl_session_cache shared:SSL:10m;
    
    # Security headers
    add_header Strict-Transport-Security "max-age=31536000; includeSubDomains" always;
    add_header X-Frame-Options "SAMEORIGIN" always;
    add_header X-Content-Type-Options "nosniff" always;
    add_header X-XSS-Protection "1; mode=block" always;
    
    location / {
        root /usr/share/nginx/html;
        try_files $uri /index.html;
    }
    
    location /api/ {
        proxy_pass http://api-gateway:8000;
        proxy_ssl_verify off;
        proxy_set_header Host $host;
        proxy_set_header X-Real-IP $remote_addr;
        proxy_set_header X-Forwarded-For $proxy_add_x_forwarded_for;
        proxy_set_header X-Forwarded-Proto $scheme;
    }
}
\end{lstlisting}

\subsection{API Rate Limiting (Advanced)}

\begin{lstlisting}[style=pythonstyle]
# services/api-gateway/middleware/rate_limiter.py
from fastapi import Request, HTTPException
from fastapi.middleware.base import BaseHTTPMiddleware
import redis
from datetime import datetime

class RateLimiter(BaseHTTPMiddleware):
    def __init__(self, app, redis_client: redis.Redis):
        super().__init__(app)
        self.redis = redis_client
    
    async def dispatch(self, request: Request, call_next):
        # Get user identifier (IP or user ID)
        user_id = request.state.user.id if hasattr(request.state, 'user') else request.client.host
        
        # Rate limit key
        key = f"rate_limit:{user_id}:{datetime.now().minute}"
        
        # Increment counter
        count = self.redis.incr(key)
        
        # Set expiry on first request
        if count == 1:
            self.redis.expire(key, 60)  # 1 minute window
        
        # Check limit
        limit = 1000  # 1000 requests per minute
        if count > limit:
            raise HTTPException(
                status_code=429,
                detail="Rate limit exceeded",
                headers={"Retry-After": "60"}
            )
        
        # Add headers
        response = await call_next(request)
        response.headers["X-RateLimit-Limit"] = str(limit)
        response.headers["X-RateLimit-Remaining"] = str(max(0, limit - count))
        
        return response

# Usage in main.py
app.add_middleware(RateLimiter, redis_client=redis_client)
\end{lstlisting}

\section{Appendix}

\subsection{Complete Service Port Reference}

\begin{table}[h]
\centering
\caption{Complete Port Mapping Reference}
\begin{tabularx}{\textwidth}{lXll}
\toprule
\textbf{Service} & \textbf{Description} & \textbf{Internal} & \textbf{External} \\
\midrule
Frontend & React application (Nginx) & 80 & 3000 \\
API Gateway & FastAPI main API & 8000 & 8000 \\
PV Classification & Classification service & 8002 & 8012 \\
RL Intelligence & RL service & 8013 & 8013 \\
PostgreSQL & Primary database & 5432 & 5432 \\
Elasticsearch & Document storage & 9200 & 9200 \\
Elasticsearch (Transport) & Node communication & 9300 & 9300 \\
Neo4j HTTP & Graph database browser & 7474 & 7474 \\
Neo4j Bolt & Graph database protocol & 7687 & 7687 \\
Redis & Cache & 6379 & 6379 \\
Prometheus & Metrics collection & 9090 & 9090 \\
Grafana & Metrics visualization & 3000 & 3001 \\
\bottomrule
\end{tabularx}
\end{table}

\subsection{Environment Variables Reference}

\begin{table}[h]
\centering
\caption{Complete Environment Variables}
\begin{longtable}{p{0.3\textwidth}p{0.3\textwidth}p{0.35\textwidth}}
\toprule
\textbf{Variable} & \textbf{Default} & \textbf{Description} \\
\midrule
POSTGRES\_USER & postgres & PostgreSQL username \\
POSTGRES\_PASSWORD & (required) & PostgreSQL password \\
POSTGRES\_DB & document\_intelligence & Database name \\
POSTGRES\_HOST & postgres & Database host \\
POSTGRES\_PORT & 5432 & Database port \\
\addlinespace
ELASTICSEARCH\_URL & http://elasticsearch:9200 & ES connection URL \\
ELASTIC\_PASSWORD & (required) & ES password \\
\addlinespace
NEO4J\_AUTH & neo4j/password & Neo4j credentials \\
NEO4J\_URL & bolt://neo4j:7687 & Neo4j Bolt URL \\
\addlinespace
REDIS\_HOST & redis & Redis host \\
REDIS\_PORT & 6379 & Redis port \\
\addlinespace
JWT\_SECRET\_KEY & (required) & JWT signing secret \\
JWT\_ALGORITHM & HS256 & JWT algorithm \\
JWT\_EXPIRATION & 1800 & Token expiry (seconds) \\
\addlinespace
OPENAI\_API\_KEY & (optional) & OpenAI API key \\
OPENAI\_MODEL & gpt-4o-mini & OpenAI model name \\
\addlinespace
REACT\_APP\_API\_URL & http://localhost:8000 & Frontend API URL \\
\bottomrule
\end{longtable}
\end{table}

\subsection{Glossary of Terms}

\begin{description}
    \item[ICSR] Individual Case Safety Report - A detailed report of a suspected adverse reaction to a drug.
    \item[Pharmacovigilance (PV)] The science of detecting, assessing, understanding, and preventing adverse effects of medications.
    \item[Ensemble Classification] Machine learning technique combining multiple models for improved accuracy.
    \item[Zero-Shot Learning] ML approach where model classifies data without task-specific training.
    \item[MeSH Terms] Medical Subject Headings - Standardized medical vocabulary by NLM.
    \item[E2B] Electronic transmission of Individual Case Safety Reports (ICH standard).
    \item[NER] Named Entity Recognition - Identifying entities (drugs, diseases) in text.
    \item[Reinforcement Learning (RL)] ML paradigm where agent learns through trial and error with rewards.
    \item[Knowledge Graph] Network representation of entities and their relationships.
    \item[WebGL] Web Graphics Library - JavaScript API for rendering 3D graphics in browsers.
    \item[JWT] JSON Web Token - Secure authentication token format.
    \item[RBAC] Role-Based Access Control - Permission system based on user roles.
\end{description}

\end{document}

